\section{Build instructions}
\label{sec:build}

%% [DATA REQUIRED]: these steps are drafted from the design files
%% (Section~\ref{sec:design-files}) and standard PCB assembly practice;
%% please verify/expand against how the physical units used in this paper
%% were actually built before submission.

\subsection{PCB fabrication and assembly}
\begin{enumerate}
\item Fabricate the \SI{45.0}{}$\times$\SI{60.0}{\milli\metre}, 2-layer
  board from \texttt{Hardware/Gerber\_PCB.zip} at any standard PCB
  manufacturer.
\item Component footprints (Table~\ref{tab:bom}) are deliberately
  hand-solder-friendly 0603 passives, except the AD5941 AFE (U1,
  \texttt{CP-48-4\_ADI-L}, a 48-pin QFN/LGA package). Two build paths:
  \begin{itemize}
  \item \textbf{Professional SMT assembly (recommended)}: submit the
    Gerbers together with \texttt{Production/\allowbreak bom-1.csv} and
    \texttt{Production/\allowbreak positions.csv} to a PCB assembly
    service; this covers U1 reliably and is the path
    \texttt{Hardware/README.md} documents for users less experienced
    with soldering.
  \item \textbf{Hand assembly}: hand-solder the 0603 passives, the
    calibration resistor bank (RCAL1/RCAL2/RCAL3), the ferrite beads, and
    the pin headers with a fine-tip iron; reflow or hot-air solder U1
    (AD5941) using its datasheet-recommended profile, since its
    fine-pitch QFN/LGA package is not practically hand-solderable with an
    iron alone.
  \end{itemize}
\item Solder or socket the Seeed Studio XIAO RP2040 (U2, DIP/pin-header
  variant) onto its footprint; this part is through-hole and
  straightforward to hand-solder.
\item Fit the appropriate calibration jumper/resistor for the intended
  measurement range if a variant other than the as-shipped
  \texttt{RCAL1}=\SI{200}{\ohm} is required (Table~\ref{tab:bom}), and
  confirm the value against \texttt{DEFAULT\_RCAL} in
  \texttt{src/data\_storage/data\_storage.h} before flashing --
  Section~\ref{sec:otherdefects}'s Defect~D-RCAL is exactly what happens
  when these two disagree.
\end{enumerate}

\subsection{Firmware build and flash}
\begin{enumerate}
\item Install the Arduino \texttt{rp2040:rp2040} core (this paper used
  version 5.6.0) and the bundled libraries under
  \texttt{Software/Arduino\_Libraries/} (Adafruit NeoPixel, an
  \texttt{arduino-printf}-based logging helper).
\item Connect the assembled board via USB; it enumerates as a standard
  RP2040 board.
\item Build and flash \texttt{AD5941\_25/\allowbreak AD5941\_25.ino}
  against the \texttt{rp2040:\allowbreak rp2040:\allowbreak
  seeed\_xiao\_rp2040} board target, e.g. with \texttt{arduino-cli}:
\begin{verbatim}
arduino-cli compile --upload -p <PORT> \
  --fqbn rp2040:rp2040:seeed_xiao_rp2040 \
  --libraries Arduino_Libraries AD5941_25
\end{verbatim}
\item Confirm a successful flash by opening a serial terminal at
  \SI{1}{\mega baud} and sending \texttt{?}; the instrument echoes its
  current parameter state (Section~\ref{sec:operation}) ending in the
  command summary line
  \texttt{E=eisScan\ \ P=SeeedStat\ \ A=CA\ \ W=SWV\ \ D=DPV\ \ M=CV\ \ O=OCP\ \ ?=params}.
\end{enumerate}
